\documentclass{article}
\usepackage{amsmath, amssymb, amsthm}
\usepackage{hyperref}
\usepackage{geometry}
\geometry{margin=1in}

\title{Erd\H{o}s Problem \#979}
\date{Last updated: 2025-08-31}
\author{Source: erdosproblems.com}

\begin{document}
\maketitle

\noindent\textbf{Status:} OPEN \quad \textbf{No prize}

\noindent\textbf{Tags:} number theory

\medskip

\noindent\textbf{Problem Statement:}

Let $k\geq 2$, and let $f_k(n)$ count the number of solutions to\[n=p_1^k+\cdots+p_k^k,\]where the $p_i$ are prime numbers. Is it true that $\limsup f_k(n)=\infty$?




Erd\H{o}s \cite{Er37b} proved this is true when $k=2$, and also when $k=3$ (but this proof appears to be unpublished).




References

[Er37b] Erd\H{o}s, Paul, On the Sum and {D}ifference of Squares of {P}rimes . J. London Math. Soc. (1937), 133--136.

\noindent\textbf{Related OEIS sequences:} A385316, possible


\bigskip
\noindent\small{Source: \url{https://www.erdosproblems.com/979}}
\end{document}
